\documentclass[10.5pt,a4paper]{article}
\usepackage[a4paper,margin=2.2cm,top=2.4cm,bottom=2.4cm]{geometry}
\usepackage{fontspec}
\usepackage{amsmath,amssymb,mathtools,bm}
\setmainfont{Noto Serif CJK SC}[BoldFont={Noto Serif CJK SC Bold}]
\setsansfont{Noto Sans CJK SC}[BoldFont={Noto Sans CJK SC Bold}]
\XeTeXlinebreaklocale "zh"
\XeTeXlinebreakskip = 0pt plus 1pt minus 0.1pt
\linespread{1.18}
\usepackage{xcolor}
\definecolor{accent}{HTML}{1F5FA8}
\definecolor{quiet}{HTML}{5A5A5A}
\definecolor{tint}{HTML}{EEF3FA}
\usepackage{titlesec}
\titleformat{\section}{\sffamily\bfseries\Large\color{accent}}{\thesection}{0.8em}{}
\titleformat{\subsection}{\sffamily\bfseries\large}{\thesubsection}{0.6em}{}
\titleformat{\subsubsection}{\sffamily\bfseries\normalsize}{\thesubsubsection}{0.5em}{}
\titlespacing*{\section}{0pt}{1.6em}{0.7em}
\titlespacing*{\subsection}{0pt}{1.2em}{0.5em}
\usepackage{enumitem}
\setlist{itemsep=0.25em,topsep=0.35em,leftmargin=1.6em}
\usepackage{array,booktabs,tabularx,xltabular,makecell}
\renewcommand{\arraystretch}{1.25}
\newcolumntype{L}{>{\raggedright\arraybackslash}X}
\newcolumntype{P}[1]{>{\raggedright\arraybackslash}p{#1}}
\usepackage{caption}
\captionsetup{font=small,labelfont=bf,labelsep=quad}
\renewcommand{\figurename}{图}
\renewcommand{\tablename}{表}
\renewcommand{\contentsname}{目录}
\usepackage{graphicx}\usepackage{float}
\usepackage{tikz}
\usetikzlibrary{positioning,arrows.meta,calc,fit}
\usepackage{fancyhdr}
\pagestyle{fancy}\fancyhf{}
\fancyhead[L]{\small\color{quiet}椭圆通道：AFRC$_{\mathrm E}$ 构件移植与遗留问题}
\fancyhead[R]{\small\color{quiet}2026 年 10 月 16 日}
\fancyfoot[C]{\small\thepage}
\renewcommand{\headrulewidth}{0.4pt}
\usepackage[colorlinks=true,linkcolor=accent,citecolor=accent,urlcolor=accent]{hyperref}
\newcommand{\AFH}{AFRC$_{\mathrm H}$}
\newcommand{\AFE}{AFRC$_{\mathrm E}$}
\newcommand{\lnJ}{$\ln J$}
\newcommand{\keyp}[1]{\textbf{\sffamily #1}}
\newenvironment{finding}{\begin{center}\begin{minipage}{0.94\linewidth}\setlength{\parskip}{0.3em}\color{black}}{\end{minipage}\end{center}}
\setlength{\parskip}{0.35em}
\setlength{\parindent}{0pt}
\tolerance=2000 \emergencystretch=3em
\usepackage{amsthm}
\theoremstyle{definition}
\newtheorem{definition}{定义}[section]
\theoremstyle{plain}
\newtheorem{proposition}[definition]{命题}
\newtheorem{corollary}[definition]{推论}
\theoremstyle{remark}
\newtheorem{remark}[definition]{注}
\renewcommand{\proofname}{证明}
\newcommand{\status}[1]{\hfill{\small\color{accent}\textsf{[#1]}}}
\newcommand{\re}{\operatorname{Re}}
\newcommand{\im}{\operatorname{Im}}
